%%%%%%%%%%%%%%%%%%%%%%%%%%%%%%%%%%%%%%%%%%%%%%%%%%%%%%%%%%%%%%%%%%%
% Who Wrote This? A Validity Audit of AI-Authorship Attribution
% in Open-Source Software Artifacts
%
% Empirical Software Engineering, Special Issue
% "Agentic Software Engineering: The Rise of AI Teammates"
%
% Build:  latexmk -pdf main.tex
%%%%%%%%%%%%%%%%%%%%%%%%%%%%%%%%%%%%%%%%%%%%%%%%%%%%%%%%%%%%%%%%%%%
\RequirePackage{fix-cm}
% natbib is a class option rather than a package load for svjour3; it is what binds
% \citep/\citet to the Springer spbasic author-year style.
\documentclass[smallextended,natbib]{svjour3}

\smartqed

\usepackage{graphicx}
\usepackage{booktabs}
\usepackage{amsmath}
\usepackage{url}
\usepackage[hidelinks]{hyperref}

% Corpus counts, generated by analysis/01_corpus/profile_aidev.py.
% Never type a corpus number into the prose: add a macro there instead.
% Generated by analysis/01_corpus/profile_aidev.py -- do not edit by hand.
% Generated 2026-08-03T14:56:39.790026+00:00
% Every corpus number in the manuscript resolves through these macros, so a
% dataset change propagates to the prose instead of silently going stale.

\newcommand{\AidevPRs}{932,791}
\newcommand{\AidevRepos}{116,423}
\newcommand{\AidevWindowStart}{2024-12-24}
\newcommand{\AidevWindowEnd}{2025-07-30}
\newcommand{\AidevTopAgentShare}{87.3\%}
\newcommand{\AidevBodyMedian}{294}
\newcommand{\AidevCommits}{88,576}
\newcommand{\AidevCommitMsgMedian}{56}
\newcommand{\AidevCommitMsgPNinety}{164}
\newcommand{\AidevTrailerRate}{39.2\%}
\newcommand{\AidevCommentBotShare}{70.1\%}
\newcommand{\AidevCurated}{33,596}
\newcommand{\AidevHumanPRs}{6,618}
\newcommand{\AidevHumanRepos}{818}
\newcommand{\AidevHumanUsers}{2,515}
\newcommand{\AidevHumanBodyMedian}{468}
\newcommand{\RepoFrame}{2,807}
\newcommand{\RepoStarsMedian}{564}
\newcommand{\LlmCutoff}{2022-11-30}
\newcommand{\AgentOpenAICodexPRs}{814,522}
\newcommand{\AgentOpenAICodexBodyMedian}{283}
\newcommand{\AgentCopilotPRs}{50,447}
\newcommand{\AgentCopilotBodyMedian}{2,681}
\newcommand{\AgentCursorPRs}{32,941}
\newcommand{\AgentCursorBodyMedian}{331}
\newcommand{\AgentDevinPRs}{29,744}
\newcommand{\AgentDevinBodyMedian}{863}
\newcommand{\AgentClaudeCodePRs}{5,137}
\newcommand{\AgentClaudeCodeBodyMedian}{1,506}

% Detector results, generated by analysis/03_validity/tables.py. Same rule.
% generated by analysis/03_validity/tables.py -- do not edit
\newcommand{\ZeroShotAurocMax}{0.498}
\newcommand{\ZeroShotAurocMin}{0.295}
\newcommand{\ZeroShotBelowChance}{11}
\newcommand{\ZeroShotCells}{11}
\newcommand{\CeilingAurocMax}{0.982}
\newcommand{\CeilingAurocMin}{0.602}
\newcommand{\CeilingRecallMax}{92.7}
\newcommand{\JudgeBotFprMax}{97.5}
\newcommand{\JudgeBotFprMin}{71.3}
\newcommand{\JudgeBotRatioMin}{16}
\newcommand{\JudgeBotRatioMax}{21}
\newcommand{\SelfAdmissionRecallMax}{3.6}
\newcommand{\SelfAdmissionFprMax}{0.0}
\newcommand{\CorrectionCells}{31}
\newcommand{\CorrectionUndefined}{8}
\newcommand{\CorrectionOutOfRange}{6}
\newcommand{\CorrectionEstimable}{17}
\newcommand{\FastDetectShortAuroc}{0.417}
\newcommand{\FastDetectLongAuroc}{0.366}
\newcommand{\CeilingShortAuroc}{0.738}
\newcommand{\CeilingLongAuroc}{0.899}
\newcommand{\PublishedBinoDiffFpr}{98.9}
\newcommand{\PublishedBinoMinFpr}{51.4}
\newcommand{\PublishedBinoMaxFpr}{98.9}
\newcommand{\PublishedFastMinFpr}{2.1}
\newcommand{\PublishedFastMaxFpr}{13.1}
\newcommand{\CeilingCorrectedTwenty}{17.1}
\newcommand{\CeilingCorrectedTwentyLo}{16.5}
\newcommand{\CeilingCorrectedTwentyHi}{17.7}

% Worked-example figures, generated by analysis/03_validity/example.py.
% generated by analysis/03_validity/example.py -- do not edit
\newcommand{\ExBinocularsScore}{-0.717}
\newcommand{\ExBinocularsThreshold}{-0.744}
\newcommand{\ExFastDetectGptScore}{1.618}
\newcommand{\ExFastDetectGptThreshold}{0.810}
\newcommand{\ExFingerprintScore}{0.853}
\newcommand{\ExFingerprintThreshold}{0.753}
\newcommand{\ExHeuristicsScore}{0.077}
\newcommand{\ExHeuristicsThreshold}{0.154}
\newcommand{\ExLlmJudgeScore}{0.100}
\newcommand{\ExLlmJudgeThreshold}{0.150}
\newcommand{\ExSelfadmissionScore}{0.000}
\newcommand{\ExSelfadmissionThreshold}{1.000}
\newcommand{\ExText}{fix: improve error handling in WebSockets}
\newcommand{\ExChars}{41}
\newcommand{\ExFlagged}{3}
\newcommand{\ExInstruments}{6}
\newcommand{\ExFlaggedNames}{Binoculars, Fast-DetectGPT and the supervised ceiling}
\newcommand{\WorkedDetector}{Binoculars}
\newcommand{\WorkedAp}{20}
\newcommand{\WorkedSe}{2.4}
\newcommand{\WorkedSp}{95.0}
\newcommand{\WorkedJ}{-0.026}
\newcommand{\WorkedNumerator}{+0.150}
\newcommand{\WorkedQuotient}{-5.9}


% Marks claims that still need a number from analysis/. Must be empty before
% submission; grep for RESULT to find every open slot.
\newcommand{\RESULT}[1]{\textbf{[#1]}}

% Marks a figure taken from someone else's paper rather than measured by us. Prints
% its argument unchanged; exists so that check_manuscript.py can tell a cited number
% from one that should have come out of analysis/.
\newcommand{\citednum}[1]{#1}

\journalname{Empirical Software Engineering}

\begin{document}

\title{Who Wrote This? A Validity Audit of AI-Authorship\\
       Attribution in Open-Source Software Artifacts}

% Without this the running head reads "Title Suppressed Due to Excessive Length".
\titlerunning{A Validity Audit of AI-Authorship Attribution in OSS}

\author{Sebastian Neumaier}

\institute{S. Neumaier \at
              University of Applied Sciences St. P\"olten,
              Campus-Platz 1, 3100 St. P\"olten, Austria \\
              \email{sebastian.neumaier@ustp.at}
           \and
              S. Neumaier \at
              Vienna University of Economics and Business,
              Welthandelsplatz 1, 1020 Vienna, Austria
}

\date{Received: date / Accepted: date}

\maketitle

\begin{abstract}
\textbf{Context.} Recent empirical work estimates how much of open-source software is
written by AI, using zero-shot detectors, keyword scans and language models as
classifiers.
\textbf{Objective.} None has been validated against ground truth on open-source
artifacts; all were benchmarked on news, essays and curated code. We ask what they
measure here, and what follows for the published estimates.
\textbf{Method.} Ground truth needs no human annotation here: agent-attributed
contributions from AIDev give a positive set of \PositiveTexts{} texts across five
genres, and artifacts from the same repositories predating ChatGPT give a negative set in
which any detection is an error by construction. We evaluate seven approaches and
propagate the measured error through the Rogan--Gladen correction.
\textbf{Results.} At the operating point published for this domain, Binoculars flags
\PublishedBinoDiffFpr{} per cent of pre-ChatGPT code diffs as AI-written, and
\PublishedBinoMinFpr{} to \PublishedBinoMaxFpr{} per cent across genres. Recalibrated to
a five per cent false positive rate the zero-shot detectors score \ZeroShotAurocMin{} to
\ZeroShotAurocMax{} AUROC, every cell below chance; against contemporaneous human text
some recovers, none rises above it. A supervised classifier on the same input reaches
\CeilingAurocMax{} on pull request bodies: the signal is present and the instruments miss
it. An instructed language model flags \JudgeBotFprMin{}--\JudgeBotFprMax{} per cent of
pre-ChatGPT bot output, responding to automation rather than to language models.
\textbf{Conclusions.} Of \CorrectionCells{} detector and genre combinations,
\CorrectionUndefined{} give an undefined correction and \CorrectionOutOfRange{} an
observed rate no prevalence could produce: most published estimates cannot be corrected
rather than merely adjusted. We give a provenance vocabulary that records an attribution
with the measured error rate of its stratum.
\keywords{AI-generated code \and Authorship attribution \and Measurement validity
\and Provenance \and Agentic software engineering}
\end{abstract}

\section{Introduction}
\label{sec:intro}

In the space of two years, a new empirical literature has formed around a single
question: how much of open-source software is now written by AI? The answers arrive as
percentages. Some fraction of commits, of pull requests, of comments, of lines of code
is reported as machine-generated, and those fractions are then used to argue about
maintainer burden, review overhead, contribution quality and the long-run
sustainability of open-source projects. The numbers circulate quickly, because the
question matters and because the answers are alarming.

Almost none of these numbers rest on a measurement instrument that has been validated
in the domain where it is being applied.

The instruments come in three families. Zero-shot statistical detectors such as
DetectGPT \citep{mitchell2023detectgpt}, Fast-DetectGPT \citep{bao2024fastdetectgpt} and
Binoculars \citep{hans2024binoculars} score text by its likelihood under a language
model. Heuristic filters match phrase lists, formatting signatures or punctuation
patterns. Large language models are prompted to classify authorship directly. Each
family was developed and benchmarked on a particular kind of text: news articles,
student essays, Wikipedia paragraphs, and code drawn from curated programming
benchmarks. Binoculars, to take the strongest published claim, reports detection of
over \citednum{90\%} of ChatGPT-generated text at a false positive rate of \citednum{0.01\%}
\citep{hans2024binoculars}. That claim was established on essay-length prose.

The artifacts of open-source software development are not essay-length prose. In the
AIDev dataset of agent-authored contributions \citep{li2025aidev}, the median commit
message is \AidevCommitMsgMedian{} characters long and the ninetieth percentile is
\AidevCommitMsgPNinety{}. The median pull request body is \AidevBodyMedian{} characters.
These texts are short, templated, written to a genre convention, and produced by a
globally distributed contributor base writing in a second language. Every one of these
properties is a known stressor for the detectors in question. Detector accuracy degrades
sharply on short inputs, and \citet{liang2023biased} showed that GPT detectors
systematically misclassify text by non-native English writers as machine-generated.
Open-source repositories additionally carry more than a decade of pre-existing
machine-generated text --- dependency bots, release automation, templated issue
reports --- that predates large language models entirely and that no published
prevalence study has been shown to exclude reliably.

If the instruments are miscalibrated in this domain, the consequences are not evenly
distributed. A detector with an inflated false positive rate does not merely add noise:
because non-AI artifacts vastly outnumber AI ones in most historical samples, even a
small false positive rate produces a large absolute number of false detections, and
prevalence is overestimated systematically rather than randomly. The same asymmetry
means that the bias documented by \citet{liang2023biased} would fall hardest on
contributors already least well served by the ecosystem.

\subsection{Auditing the instrument rather than the phenomenon}
\label{sec:intro:inversion}

This paper inverts the question. Rather than adding another prevalence estimate, it
asks what the existing measurement instruments actually measure when applied to
open-source artifacts, and what follows for the estimates already in the literature.

The inversion is made tractable by two sources of ground truth that require no human
annotation --- an essential property, since human annotation of authorship is precisely
what is unavailable at scale, and inter-rater reliability on ``does this look
AI-written'' is exactly the thing under dispute.

\paragraph{Positives.} Agentic coding platforms leave attribution traces. When Devin,
OpenAI Codex, GitHub Copilot, Cursor or Claude Code opens a pull request under its own
account, authorship is declared by the platform rather than inferred. The AIDev
dataset \citep{li2025aidev} collects \AidevPRs{} such pull requests across
\AidevRepos{} repositories. A second declared channel, the \texttt{Co-authored-by}
trailer, appears on \AidevTrailerRate{} of the commits in that dataset.

\paragraph{Negatives.} Text written before large language models were publicly
available cannot have been produced by them. Artifacts committed to the same
repositories before the release of ChatGPT on \LlmCutoff{} therefore form a negative set
whose purity follows from chronology rather than from judgement. This yields a
counterfactual test with an unusual property: every artifact a detector flags in this
set is a false positive \emph{by construction}, with no annotation, adjudication or
appeal. The resulting false positive rate is not an estimate of detector error under
assumptions --- it is a direct measurement.

Given measured sensitivity and specificity, published prevalence figures can then be
corrected. The estimator is not new; it is standard in epidemiology, where the problem
of inferring true disease prevalence from an imperfect screening test was solved by
\citet{rogan1978estimating} nearly fifty years ago. Its application here is
straightforward, and its absence from the software engineering literature on AI
prevalence is, we argue, the central methodological gap.

\subsection{Why this matters now}
\label{sec:intro:now}

The scale of the downstream literature makes the gap consequential. The 2026 Mining
Software Repositories Mining Challenge was built on AIDev and produced \citednum{62} accepted
papers, studying merge and rejection rates, test coverage, continuous integration
outcomes, security defects, refactoring behaviour, code clones and reviewer sentiment
in agent-authored contributions. All of them take the attribution labels as given ---
reasonably so, since AIDev's labels are platform-declared. But an entire research
programme has now been built on top of authorship attribution in under two years,
while the question of what attribution-by-detection actually measures, in the cases
where declaration is absent, has gone unasked.

That gap is beginning to be noticed from one side. \citet{khosravani2026census} audit
the \emph{declaration} channels --- configuration files, commit-message markers,
author identity, bot signatures --- across \citednum{180} million repositories, and find that
single-signal detection undercounts badly: bot-account lookup alone recovers only \citednum{3.3\%}
of the Claude Code commits their multi-method approach identifies. Their finding and
ours are complementary and point in opposite directions. They show that
\emph{declaration-based} measurement misses AI contributions that are real. We ask
whether \emph{content-based} measurement invents AI contributions that are not. Between
the two, a prevalence estimate can be bounded from both sides rather than reported as a
point.

\subsection{Contributions}
\label{sec:intro:contributions}

\begin{enumerate}
  \item \textbf{A criterion-validity evaluation} of the detection approaches used in the
    open-source prevalence literature, against attribution-based ground truth, reported
    per agent and per artifact genre rather than pooled (Section~\ref{sec:validity}).
    Pooling matters more than it may appear: \AidevTopAgentShare{} of AIDev is a single
    agent, so a pooled metric is that agent's metric wearing a general label.
  \item \textbf{A counterfactual false positive measurement} on pre-ChatGPT artifacts,
    where every detection is an error by construction, stratified by genre, language,
    text length and project characteristics, and separated from the confounding
    contribution of pre-LLM automation (Section~\ref{sec:validity}).
  \item \textbf{Bias-corrected prevalence estimates} that propagate the measured error
    rates through the Rogan--Gladen estimator, reported as intervals under best- and
    worst-case detector assumptions and contrasted with the point estimates in the
    literature (Section~\ref{sec:correction}).
  \item \textbf{An interoperable provenance model} that represents authorship claims
    together with their evidence type --- self-declared trailer, platform metadata,
    agent account, or detector inference with a score --- so that an attribution can be
    inspected and queried rather than taken on faith
    (Section~\ref{sec:provenance}).
  \item \textbf{A replication package} containing the corpus specification, the detector
    harness, the provenance vocabulary and all analysis code.
\end{enumerate}

The argument of this paper is not that AI-generated content in open source is rare, nor
that the concern motivating the prevalence literature is misplaced. It is that the
field currently cannot tell the difference between a detector finding AI text and a
detector finding short text, templated text, or text written by someone whose first
language is not English --- and that this is a measurable, correctable state of affairs
rather than an inherent limitation.

\section{Background and related work}
\label{sec:background}

\subsection{Prevalence studies of AI-generated contributions}
\label{sec:background:prevalence}

\RESULT{Table~1: the audit targets. One row per published prevalence study, with
columns for detection strategy, artifact types scored, sample frame, reported
prevalence, and whether in-domain validation was performed. Populate from
LITERATURE.md section A once the four 2026 studies have been read in full and their
pipelines extracted. The final column is expected to read ``none'' throughout; that
column is the paper's motivation in one glance.}

\subsection{Detecting machine-generated text}
\label{sec:background:detection}

Zero-shot detection exploits a structural asymmetry: text sampled from a language model
tends to sit at a local maximum of that model's likelihood surface, while human text
does not. \citet{mitchell2023detectgpt} formalised this as probability curvature and
estimated it by perturbing the candidate text and re-scoring, which is accurate but
computationally heavy. \citet{bao2024fastdetectgpt} replaced the perturbation step with
a conditional probability curvature computed analytically, achieving comparable accuracy
at a fraction of the cost. \citet{hans2024binoculars} took a different route, contrasting
the perplexity of the text under one model with its cross-perplexity under a second,
closely related model; the ratio normalises away the intrinsic difficulty of the text
and yields the strongest published operating points.

Two properties of this family matter here. All of it is calibrated on text of
substantial length: the quantities being estimated are sample means over tokens, and
their variance grows as texts shorten. The null hypothesis is also ``written by a
human'' rather than ``not written by a language model''. Boilerplate produced by
non-LLM automation sits in the same low-perplexity region as model output, and nothing
in the method separates the two.

Code-specific detection adapts the same intuition to a different surface form.
\RESULT{DetectCodeGPT citation and one-paragraph method summary --- verify the exact
venue, listed as unresolved in LITERATURE.md section B.}

The closest work to ours in the detection family is \citet{ghaleb2026fingerprinting},
who trains a supervised classifier on \citednum{41} engineered features spanning commit messages,
pull request structure and code characteristics, and reports \citednum{97.2\%} F1 at identifying
\emph{which} of five agents authored a pull request. The task differs from ours: it
discriminates among agents on contributions already known to be agent-authored, so it
says nothing about the agent-versus-human decision or about false positives on human
text. We include it as a supervised, in-domain reference point. It estimates what is
achievable when a classifier is trained on the distribution it is tested on, and gives a
scale against which the off-the-shelf detectors can be read.

\subsection{Detector validity and bias}
\label{sec:background:validity}

The reliability of AI-text detection has been contested since the tools appeared.
\citet{sadasivan2023reliably} establish theoretical limits, showing that as the
distributions of human and machine text converge, any detector's performance approaches
chance, and that paraphrase attacks exploit this directly.

The empirical result most relevant to our setting is \citet{liang2023biased}, who found
that detectors classify writing by non-native English speakers as machine-generated at
substantially elevated rates, because the linguistic features that detectors treat as
evidence of machine authorship --- limited lexical variety, conventional phrasing,
lower syntactic complexity --- also characterise second-language writing. Open-source
development is a setting where this bias would be both amplified and consequential:
contributors are globally distributed, most write in English as a second language, and
being flagged as an AI contributor carries a governance cost.

Validation of detectors outside their original domain remains rare. In the news domain,
\citet{ansari2025echoes} apply Binoculars, Fast-DetectGPT and GPTZero to over
\citednum{40,000} articles to estimate the growth of LLM use in journalism --- the same instrument stack
as the software prevalence literature, applied with the same absence of in-domain
validation. \RESULT{If the pre-GPT-era Washington Post counterfactual baseline is
confirmed (see LITERATURE.md section C, currently unverified), cite it here as the
direct methodological precedent for our design and state explicitly what our study adds:
the software domain, code as well as prose, and the correction step.}

\subsection{Prevalence estimation under imperfect measurement}
\label{sec:background:correction}

Estimating the true frequency of a condition from an imperfect test is a solved problem
outside software engineering. \citet{rogan1978estimating} give the standard correction:
for an apparent prevalence $p_{\text{obs}}$ measured by a test with sensitivity $s$ and
specificity $c$, the true prevalence is
\begin{equation}
  p_{\text{true}} = \frac{p_{\text{obs}} + c - 1}{s + c - 1}.
  \label{eq:rogan-gladen}
\end{equation}
The estimator makes the dependence explicit. When $s + c - 1$ is small, meaning the test
barely outperforms chance, the correction amplifies uncertainty without bound and the
corrected estimate carries no information. Applying it to published figures needs the
quantities that have not been measured for this domain, which
Section~\ref{sec:validity} supplies.

\subsection{Bots and pre-LLM automation}
\label{sec:background:bots}

Machine-generated text in open-source repositories predates large language models by a
decade. Dependency managers, release automation, continuous integration reporters and
issue templates all produce text at scale, and identifying them is an established
research problem \RESULT{BoDeGHa citation, verify}. The confounder is measurable: in
AIDev, \AidevCommentBotShare{} of comments on agent-authored pull requests come from
bots. An estimate that does not separate bot text from model text measures their sum,
and a detector evaluation that does not report bot text separately cannot say which of
the two it responds to.

\subsection{Provenance and attribution infrastructure}
\label{sec:background:provenance}

\RESULT{PROV-O \citep{lebo2013provo} as the base vocabulary; SLSA and in-toto as the
industrial supply-chain counterpart; existing software-provenance modelling work. Frame
the gap: supply-chain provenance answers ``what produced this artifact'' for build
systems, but has no vocabulary for ``who or what authored this content, and on what
evidence''. Draft once the SLSA/in-toto positioning search in LITERATURE.md section F
is complete.}

\section{Ground-truth construction}
\label{sec:groundtruth}

The validity of everything that follows depends on this section, so we describe the
corpus in more detail than is conventional. Two design commitments govern it. First,
no label originates from human judgement about whether a text looks machine-written;
every label derives from platform attribution or from chronology. Second, positives and
negatives are drawn from the same repositories wherever possible, so that differences
between them are not differences between projects.

\subsection{Positive set: declared agent authorship}
\label{sec:groundtruth:positives}

Positives come from the AIDev dataset \citep{li2025aidev}, which collects pull requests
opened by agentic coding platforms operating under their own accounts. We pin the
dataset revision used, because AIDev is a living resource: the paper introducing it
reports \citednum{456{,}000} pull requests, while the revision we analyse contains \AidevPRs{}
across \AidevRepos{} repositories. All counts below are measured from the pinned
revision rather than quoted from the literature.

Three properties of the positive set constrain the design, and we report them
prominently because they bound every claim in this paper.

\paragraph{The observation window is seven months.} All agent-authored pull requests in
the dataset fall between \AidevWindowStart{} and \AidevWindowEnd{}. Our findings
therefore concern the agent platforms and underlying model versions current in the first
half of 2025. We return to this in Section~\ref{sec:threats}.

\paragraph{The agent mix is severely imbalanced.} Table~\ref{tab:agents} gives the
composition. A single agent accounts for \AidevTopAgentShare{} of all pull requests, and
median body length differs by roughly an order of magnitude across agents. Any metric
computed by pooling agents is therefore dominated by one of them. We report all
detector results per agent, and present pooled figures only under explicit equal
allocation.

\begin{table}
\caption{Composition of the positive set at the pinned AIDev revision. Median body
length varies by roughly an order of magnitude across agents, so pooled metrics are
dominated by the largest contributor.}
\label{tab:agents}
\begin{tabular}{lrr}
\hline\noalign{\smallskip}
Agent & Pull requests & Median body (chars) \\
\noalign{\smallskip}\hline\noalign{\smallskip}
OpenAI Codex & \AgentOpenAICodexPRs{} & \AgentOpenAICodexBodyMedian{} \\
GitHub Copilot & \AgentCopilotPRs{} & \AgentCopilotBodyMedian{} \\
Cursor & \AgentCursorPRs{} & \AgentCursorBodyMedian{} \\
Devin & \AgentDevinPRs{} & \AgentDevinBodyMedian{} \\
Claude Code & \AgentClaudeCodePRs{} & \AgentClaudeCodeBodyMedian{} \\
\noalign{\smallskip}\hline
\end{tabular}
\end{table}

\paragraph{The texts are short.} The median pull request body is \AidevBodyMedian{}
characters and the median commit message \AidevCommitMsgMedian{}. This is not incidental
to the study; it is arguably the central fact about the domain. Detectors that report
near-perfect separation on essay-length text are being asked here to decide on a single
line of English.

\subsubsection{Artifact genres}
\label{sec:groundtruth:genres}

We treat five genres separately, because they differ in length, register and
conventionality, and because the prevalence literature scores different subsets of them:

\begin{enumerate}
  \item \emph{Pull request descriptions} --- prose, moderate length, often templated.
  \item \emph{Commit messages} --- the shortest genre, heavily conventionalised.
    \texttt{Co-authored-by} and similar trailer lines are stripped into a separate field
    and never passed to a detector; recovering an explicit declaration is not detection.
  \item \emph{Code diffs} --- the patch hunks, scored per file, with language taken from
    repository metadata.
  \item \emph{Issue bodies} --- prose with a different register from pull requests.
  \item \emph{Pull request comments} --- restricted to human accounts. This restriction
    is severe and necessary: \AidevCommentBotShare{} of comments on agent-authored pull
    requests are bot-authored, so an unfiltered comment genre would measure
    continuous-integration bot output rather than agent output.
\end{enumerate}

Allocation across the resulting agent-by-genre cells is equal rather than proportional,
targeting 1{,}000 texts per cell with a floor of 300. Equal allocation sacrifices
precision on the most abundant agent to obtain usable precision on the least abundant,
which is the correct trade when the question is whether detector behaviour generalises
across agents rather than what its average is. The sampled positive set contains
\PositiveTexts{} texts across \PositiveCells{} cells.

\paragraph{What equal allocation costs.} Because agents differ so much in verbosity,
reweighting them equally changes the corpus itself: the median pull request body rises
from \PrBodyMedianPopulation{} characters in the population to
\PrBodyMedianSampled{} in the sample. The sampled corpus is therefore deliberately
\emph{not} representative of the population of agent contributions, and no prevalence
claim is made from it --- it exists to estimate per-stratum detector error, which
Section~\ref{sec:correction} then applies to representative samples. Reporting this
explicitly matters because the same reweighting, left undisclosed, would make detectors
look better than they are: longer texts are easier to classify.

\paragraph{One genre cannot support per-agent analysis.} Issue bodies are linked to
agent authorship only through AIDev's issue--pull-request link table, and the resulting
cells are thin: three of the five agents fall below the floor of 300, two of them below
200. Issue bodies are therefore analysed pooled across agents, and per-agent claims are
restricted to the other four genres. This is a limitation of the available attribution
data rather than of the design, and we report it rather than quietly padding the cells
from a proportional draw.

\subsection{Negative sets}
\label{sec:groundtruth:negatives}

No single negative set is simultaneously pure and representative, so we construct three
with different and complementary weaknesses.

\subsubsection{N1: pre-LLM artifacts from the same repositories}
\label{sec:groundtruth:n1}

The primary negative set consists of artifacts written in the positive set's own
repositories before \LlmCutoff{}, the public release of ChatGPT. Its purity follows from
chronology, and every detection in it is a false positive by construction.

\paragraph{Dating the text, not the artifact.} Filtering on an artifact's creation date
is insufficient, because bodies can be edited and issue threads accumulate comments
indefinitely. While profiling repository data we encountered an issue created in 2013
carrying a comment written in 2026 that opens with an explicit agent-generated
verification report. Selecting on creation date alone would have admitted that comment
to the pre-LLM negative set. We therefore require both the creation and last-modification
timestamps to precede the cutoff, which guarantees the text has not been touched since,
at the cost of biasing the set toward artifacts that attracted no subsequent activity.
Section~\ref{sec:threats} treats this quiescence bias.

For commit messages and diffs the problem does not arise in the same form: we read them
from cloned repository history rather than from the platform API, where they are dated
by committer timestamp and cannot be altered without rewriting history.

\paragraph{Repository eligibility.} Not every repository in the positive set has a
pre-LLM history. Of the \RepoFrame{} curated repositories, \NOneResolved{} still resolve
on GitHub --- \NOneUnresolvable{} have been deleted, renamed or withdrawn --- and of
those, \NOneEligible{} (\NOneEligiblePct{}) were created before \LlmCutoff{}. These form
the N1 frame. The oldest dates from \NOneOldestRepo{}, so the frame spans the better part
of two decades of development practice.

Repository names are not stable identifiers: a project can be renamed and its name
reused by an unrelated project. We therefore verify that each name still resolves to the
GitHub numeric identifier that AIDev recorded, and exclude any repository where it does
not, independently of age.

The excluded \NOneExcluded{} repositories are not a random subset --- they are younger, and younger projects plausibly differ in their
propensity to adopt coding agents. Losing nearly two-fifths of the frame is a real cost
to statistical power in the matched analysis, and we report a generalisation check on
the excluded repositories using positives alone to test whether the restriction changes
the conclusions.

\subsubsection{N2: contemporaneous human contributions}
\label{sec:groundtruth:n2}

N1 is pure but temporally separated from the positives, which confounds authorship with
period: coding conventions, tooling and project norms all changed between 2022 and 2025.
The second negative set removes the temporal confound at the cost of purity. AIDev
includes \AidevHumanPRs{} human-authored pull requests from \AidevHumanRepos{}
repositories and \AidevHumanUsers{} distinct users, drawn from the same period as the
positives, with a median body of \AidevHumanBodyMedian{} characters.

These are human-\emph{attributed}, not verified human-\emph{written}: a developer may
have drafted the text with a chatbot and pasted it in, leaving no trace. Error rates on
N2 are therefore reported as \emph{apparent} rather than true false positive rates. The
gap between the N1 and N2 rates is itself informative, since under stated assumptions it
bounds the extent of undeclared assisted authorship in the human population --- the
stratum this study otherwise cannot observe.

\subsubsection{N3: pre-LLM automation}
\label{sec:groundtruth:n3}

The third set isolates the confounder. It comprises pre-cutoff artifacts from known
automation accounts --- dependency managers, release tooling, continuous integration
reporters --- in the same repositories. This text is machine-generated and definitively
not LLM-generated. The detection rate on N3 estimates how much reported ``AI
prevalence'' is decade-old automation being relabelled.

\subsection{Matching and stratification}
\label{sec:groundtruth:matching}

Because detector scores are length-sensitive and the agents differ substantially in
length, an unmatched comparison would partly measure length rather than authorship.
Negatives are therefore sampled to match the positive length distribution within genre
by decile. We report unmatched results alongside matched ones, since the difference
between them estimates the magnitude of the length confound rather than merely
controlling for it.

Stratification variables are programming language, text length, artifact genre, agent,
and project scale. For the subset of repositories also covered by the CrOSSD panel
\RESULT{cite Dam, Klausner, Neumaier WWW'23} we additionally use its health and
criticality metrics. Following \citet{liang2023biased}, we require a proxy for
contributors' English proficiency, but author-name and author-location proxies are both
unreliable and ethically objectionable. We instead use repository-level linguistic
proxies derived from documentation text, which attribute a property to the project
rather than to an individual.

\subsection{Corpus release}
\label{sec:groundtruth:release}

The frozen corpus is released as a specification --- artifact identifiers, selection
predicates, seeds and per-cell counts --- together with the code that rebuilds it, so
that the corpus can be reconstructed without redistributing repository content beyond
what its licence permits. Every count reported in this paper is emitted directly by the
analysis scripts into the manuscript source, so no figure in the prose can drift from the
data that produced it.

\section{Detector validity (RQ1 and RQ2)}
\label{sec:validity}

\subsection{Detector harness}
\label{sec:validity:harness}

Every approach is wrapped in one interface: a list of texts goes in, a list of scores
comes out, higher meaning more likely machine-written. Adapters normalise the direction
of statistics that natively point the other way. Nothing else reaches an adapter. The
driver passes strings rather than a data frame, so no agent label, repository, date or
set membership is visible to a detector, and scores are joined back to their rows
afterwards. The guard matters because most fields in the corpus correlate with the
label: a detector shown a 2019 timestamp would not need to read the text.

Seven instruments run, in four groups. The zero-shot statistical detectors are
Fast-DetectGPT \citep{bao2024fastdetectgpt}, Binoculars \citep{hans2024binoculars} and
DetectCodeGPT \citep{shi2025betweenlines}, the stack \citet{ji2026exploratory} apply.
The keyword group is the self-admission scan of \citet{mao2026measurement} and
\citet{almujahid2026aiusage}, reconstructed from their described list of six prefix words
crossed with eleven tool names. A set of formatting and phrasing rules stands for the
screening maintainers apply informally; it reproduces no published pipeline and is
reported as its own rule set rather than as anyone's method. An instructed language
model is the fourth, prompted once with a prompt fixed before any result was seen. The
supervised classifier of Section~\ref{sec:validity:ceiling} is the reference point
rather than a candidate instrument.

Scoring models are pinned with their weights recorded in each adapter's version string.
Fast-DetectGPT scores under Qwen2.5-7B, Binoculars under the falcon-7b pair of its
original publication, and DetectCodeGPT under CodeLlama-7b, which is the base model of
its reference implementation and the configuration \citet{ji2026exploratory} deploy. No
commercial detector is included, so the evaluation reproduces from open weights.

Reimplementing DetectCodeGPT exposed a discrepancy worth recording, since it bears on
reproducibility in this literature generally. The paper's Algorithm~1 subtracts a mean
of normalised perturbed log rank terms, which contradicts its own definition of the
statistic and the released code; and the paper gives a Poisson mean of two for newline
insertion where the released code inserts exactly one, its Poisson draw commented out.
We follow the released code, on the grounds that it is the artifact that produced the
published numbers, and record the choice in the adapter version string.

\subsection{RQ1: criterion validity}
\label{sec:validity:rq1}

Table~\ref{tab:rq1} reports discrimination against the length-matched pre-ChatGPT
negatives drawn from the same repositories.

% generated by analysis/03_validity/tables.py -- do not edit
\begin{table}[t]
\caption{Discrimination against length-matched pre-ChatGPT negatives from the same
repositories. Above: area under the ROC curve, where 0.5 is chance. Below: recall at the
threshold admitting a 5\% false positive rate on those negatives.}
\label{tab:rq1}
\small
\begin{tabular}{@{}lrrrrr@{}}
\toprule
& \multicolumn{5}{c}{Artifact genre} \\
\cmidrule(l){2-6}
Instrument & Commit msg & PR body & Issue & Comment & Diff \\
\midrule
\multicolumn{6}{@{}l}{\emph{AUROC}} \\
Binoculars & 0.471 & 0.467 & 0.454 & 0.413 & 0.459 \\
Fast-DetectGPT & 0.498 & 0.295 & 0.410 & 0.390 & 0.455 \\
Formatting rules & 0.575 & 0.711 & 0.639 & 0.560 & 0.534 \\
Self-admission & 0.503 & 0.518 & 0.505 & 0.503 & 0.502 \\
LLM judge & 0.635 & 0.939 & 0.631 & 0.619 & 0.638 \\
Supervised (ceiling) & 0.864 & 0.982 & 0.877 & 0.769 & 0.602 \\
\addlinespace
\multicolumn{6}{@{}l}{\emph{Recall at 5\% false positive rate (\%)}} \\
Binoculars & 2.4 & 6.0 & 2.4 & 1.4 & 2.5 \\
Fast-DetectGPT & 5.7 & 0.3 & 3.0 & 2.5 & 5.3 \\
Formatting rules & 12.3 & 15.4 & 11.7 & 8.3 & 3.1 \\
Self-admission & 0.6 & 3.6 & 0.9 & 0.6 & 0.3 \\
LLM judge & 20.7 & 80.0 & 29.7 & 27.8 & 14.1 \\
Supervised (ceiling) & 46.0 & 92.7 & 55.6 & 29.1 & 16.9 \\
\bottomrule
\end{tabular}
\end{table}


The zero-shot detectors do not separate the classes. Across all \ZeroShotCells{}
combinations of detector and genre, Binoculars and Fast-DetectGPT score between
\ZeroShotAurocMin{} and \ZeroShotAurocMax{}, and every one of them falls below 0.5.
Below chance is not a weaker version of the same finding. It means that in this corpus
the two instruments rank human text written before ChatGPT existed as more
machine-written than text an agent actually wrote, consistently and in every genre.
Fast-DetectGPT on pull request bodies, at \ZeroShotAurocMin{}, would be a usable
detector if its decision were reversed.

Recall at a fixed operating point makes the practical consequence clearer than the
AUROC. Calibrated to flag five per cent of pre-ChatGPT text, Binoculars recovers between
one and six per cent of agent-authored artifacts, which is at or below what flagging at
random would return. The prevalence literature applies these instruments at published
thresholds rather than at thresholds fitted in domain, and
\citet{ji2026exploratory} report the thresholds themselves moving by a factor of four
across the domains they surveyed.

The self-admission scan behaves in the opposite way and exactly as its mechanism
predicts. Its false positive rate on pre-ChatGPT text is \SelfAdmissionFprMax{} per cent
in every genre, since the tool names it searches for did not exist, and its recall never
exceeds \SelfAdmissionRecallMax{} per cent, since an agent has no reason to write
``generated by ChatGPT'' into its own commit message. It is a precise instrument for a
quantity close to zero. An estimate built on it is a lower bound whose distance from the
truth is unknown, which is the complement of the result \citet{khosravani2026census}
report for the metadata channels.

\subsubsection{The supervised ceiling}
\label{sec:validity:ceiling}

A detector can fail because the signal is absent or because the instrument misses it.
Separating the two needs a classifier that is allowed to see the training distribution.
Following \citet{ghaleb2026fingerprinting}, we fit gradient-boosted trees on features
computed from the artifact text, cross-validated in folds grouped by repository so that
no project's house style is learned in training and rewarded in test. His temporal
features are excluded rather than unavailable: submission hour and day of week separate
our classes perfectly and for the wrong reason, since the negatives predate
2022-11-30 by construction. The remaining 24 features are text-derived, which is the
point rather than a limitation, because a ceiling that sees more than the detectors it
bounds is not a ceiling.

The ceiling ranges from \CeilingAurocMin{} on diffs to \CeilingAurocMax{} on
\CeilingBestGenre{}. The signal is therefore present in most genres and substantial in
one, and the zero-shot detectors are not finding it. Diffs are the exception in both
directions: the supervised model reaches only \CeilingAurocMin{} there, so the failure
on code is partly a property of the artifact rather than of the instruments.

\subsection{RQ2: structure of the errors}
\label{sec:validity:rq2}

\subsubsection{The counterfactual false positive test}
\label{sec:validity:counterfactual}

Every detection in the pre-ChatGPT set is an error by construction, so the rate needs no
assumption about what the negatives contain. The rates in Table~\ref{tab:rq1} are held
at five per cent by calibration, which makes the recall column the informative one; the
question this subsection asks is which pre-ChatGPT text gets flagged, and Table~\ref{tab:rq2automation}
answers it.

\subsubsection{Automation as a confounder}
\label{sec:validity:bots}

% generated by analysis/03_validity/tables.py -- do not edit
\begin{table}[t]
\caption{False positives on pre-ChatGPT automation (\%), at the threshold calibrated to
flag 5\% of pre-ChatGPT human text in the same genre. Both sets predate language models,
so every detection in either is an error; the difference is what the instrument responds
to. Diffs carry no bot stratum.}
\label{tab:rq2automation}
\small
\begin{tabular}{@{}lrrrrr@{}}
\toprule
Instrument & Commit msg & PR body & Issue & Comment & Diff \\
\midrule
Binoculars & 16.7 & 1.6 & 0.7 & 12.7 & --- \\
Fast-DetectGPT & 7.5 & 16.1 & 2.2 & 4.5 & --- \\
Formatting rules & 8.6 & 0.3 & 2.1 & 11.3 & --- \\
Self-admission & 0.0 & 0.0 & 0.0 & 0.0 & --- \\
LLM judge & 72.3 & 97.5 & 80.6 & 71.3 & --- \\
Supervised (ceiling) & --- & --- & --- & --- & --- \\
\bottomrule
\end{tabular}
\end{table}


Machine-generated text in repositories predates language models by a decade, and the
instructed language model responds to it far more strongly than to the human text of the
same period. Held at a five per cent false positive rate on pre-ChatGPT human writing,
it flags between \JudgeBotFprMin{} and \JudgeBotFprMax{} per cent of pre-ChatGPT bot
output, a factor of \JudgeBotRatioMin{} to \JudgeBotRatioMax{}. Dependency updates,
release notes and continuous integration reports were written by scripts, and the judge
is answering a question about machine authorship in general rather than about language
models in particular. A prevalence estimate produced by such a judge counts a decade of
automation as AI adoption.

The formatting rules show the same pattern more weakly, which is expected of rules that
match on structure that templates also produce. Binoculars and Fast-DetectGPT do not
show it consistently, and on issue bodies and pull request bodies they flag bot text
less often than human text.

\subsubsection{Text length}
\label{sec:validity:length}

Detector statistics are sample means over tokens, so their variance falls as texts
lengthen, and performance should improve with length. It does for the supervised
ceiling, which rises from \CeilingShortAuroc{} in the shortest quintile to
\CeilingLongAuroc{} in the longest, and for the instructed model and the formatting
rules. Fast-DetectGPT moves the other way, from \FastDetectShortAuroc{} to
\FastDetectLongAuroc{}. Length is therefore not a sufficient explanation for the
zero-shot failures: were short text the whole problem, the longest quintile would be
where these detectors recover, and it is where one of them is furthest below chance.

\subsubsection{Agent generalisation}
\label{sec:validity:agents}

Because \AidevTopAgentShare{} of AIDev is a single agent, the corpus allocates the
positive set equally across the five and every metric above is reported per genre with
that allocation. Per-agent results are in the replication package. The variation across
agents is smaller than the variation across genres for every instrument, so the failures
reported here are not an artifact of one platform's house style.

\subsection{Corrected estimates (RQ3)}
\label{sec:correction}

Equation~\eqref{eq:rogan-gladen} recovers prevalence from an observed rate given
sensitivity and specificity. Its denominator is Youden's $J = s + c - 1$, reported in
Table~\ref{tab:rq3}.

% generated by analysis/03_validity/tables.py -- do not edit
\begin{table}[t]
\caption{Youden's $J = s + c - 1$, the denominator of the Rogan--Gladen correction, at
the 5\% operating point. The correction divides the observed rate by this quantity, so a
value near zero means the corrected estimate carries no information whatever the
observed rate.}
\label{tab:rq3}
\small
\begin{tabular}{@{}lrrrrr@{}}
\toprule
Instrument & Commit msg & PR body & Issue & Comment & Diff \\
\midrule
Binoculars & -0.026 & +0.010 & -0.026 & -0.036 & -0.025 \\
Fast-DetectGPT & +0.007 & -0.047 & -0.020 & -0.024 & +0.003 \\
Formatting rules & +0.102 & +0.149 & +0.080 & +0.071 & +0.023 \\
Self-admission & +0.006 & +0.036 & +0.009 & +0.006 & +0.003 \\
LLM judge & +0.162 & +0.741 & +0.258 & +0.235 & +0.105 \\
Supervised (ceiling) & +0.410 & +0.877 & +0.506 & +0.241 & +0.119 \\
\bottomrule
\end{tabular}
\end{table}


The correction cannot be applied to most of this table. Of \CorrectionCells{} detector
and genre combinations, \CorrectionUndefined{} have a bootstrap confidence interval for
$J$ that covers zero, which makes the estimator undefined rather than imprecise. In a
further \CorrectionOutOfRange{}, an observed rate of five per cent implies a corrected
prevalence above one. That is not a large estimate but a contradiction: no prevalence in
the unit interval could have produced that many flags through an instrument with that
sensitivity, so the flagged set contains something the instrument was not measuring. The
remaining \CorrectionEstimable{} cells admit a correction at an observed rate of five per
cent, and most of them stop admitting one by twenty per cent.

The consequence for the published estimates is not that they are too high by some
factor we can supply. It is that the instruments behind them do not support a corrected
estimate at all in this domain. Where a correction is available it still moves the
number: at an observed rate of twenty per cent, the supervised ceiling on pull request
bodies corrects to \CeilingCorrectedTwenty{}~per cent
(\CeilingCorrectedTwentyLo{}--\CeilingCorrectedTwentyHi{}), and the instruments actually
used in the literature give either an undefined result or a contradiction.

We make no re-estimate of prevalence on a fresh sample. Doing so would require an
instrument with a $J$ bounded away from zero on the genre being counted, and
Table~\ref{tab:rq3} does not identify one among the detectors the literature uses. The
honest output of this design is the negative result, together with the measurement that
establishes it.

\section{Recording attribution and its evidence (RQ4)}
\label{sec:provenance}

Sections~\ref{sec:validity} and~\ref{sec:correction} measure how far a detector score is
from an authorship fact. That distance is a property of the score, the genre and the
threshold, and it is known once measured. The problem is that nothing in current
practice carries it: a prevalence figure, a dashboard flag or a moderation decision
records that an artifact was classified as AI-written and discards the grounds. This
section describes a small vocabulary that keeps the two together, and shows the
questions that become answerable once it does.

\subsection{The model}
\label{sec:provenance:model}

\texttt{aiprov} extends PROV-O \citep{lebo2013provo} with one commitment: an authorship
claim is a thing in its own right, carrying its own evidence, rather than a label on an
artifact. An \texttt{aiprov:AuthorshipClaim} points at an \texttt{aiprov:Artifact} and,
where one is named, at an \texttt{aiprov:CodingAgent} identified by tool, model and
version, because the same tool across two model versions is not the same author. Each
claim carries at least one \texttt{aiprov:Evidence}.

Evidence divides by how much it rests on inference. Table~\ref{tab:evidence} lists the
classes.

\begin{table}[t]
\caption{Evidence classes in \texttt{aiprov}, ordered by how far the attribution is from
a statement by the party that made the artifact.}
\label{tab:evidence}
\small
\begin{tabular}{@{}p{0.26\textwidth}p{0.30\textwidth}p{0.34\textwidth}@{}}
\toprule
Class & Where it comes from & What it carries \\
\midrule
Agent account & The platform opened the contribution under its own account &
Agent identity; no error rate \\
\addlinespace
Trailer & A \texttt{Co-authored-by} line naming an agent &
Agent identity; no error rate \\
\addlinespace
Platform metadata & Other fields the platform records &
Agent identity; no error rate \\
\addlinespace
Self-admission & The artifact text names the tool that wrote it &
Agent identity; reliable where present, silent otherwise \\
\addlinespace
Detector & An instrument scored the text above a threshold &
Score, threshold, detector version and weights, and the stratum's false positive rate,
sensitivity and Youden's $J$ \\
\bottomrule
\end{tabular}
\end{table}

The first four rows are declared: authorship is stated rather than inferred, and there
is no error rate to record because nothing was estimated. The last row is the inferred
case, and it is the only one that carries the measurements of
Sections~\ref{sec:validity:rq1} and~\ref{sec:correction}. Attaching them to the evidence
rather than to the detector is what lets a consumer weigh one claim against another: the
same score means different things on a pull request body and on a diff, and the stratum
is what says so.

Two modelling choices are worth stating. The first is that an artifact with no claim is
not an artifact claimed to be human-written. Absence of a claim represents absence of
evidence, and conflating the two is how a detector's silence becomes a positive
assertion about a contributor. The second is that competing claims about one artifact
coexist. A pull request declared by an agent account and flagged by three detectors
carries four claims, and a consumer can ask which of them would survive on its own.

\subsection{Demonstration}
\label{sec:provenance:demo}

The corpus, its detector scores and the stratum error rates of
Tables~\ref{tab:rq1} and~\ref{tab:rq3} are loaded into a graph over a sample of the
corpus. Five competency questions are answered in SPARQL, and each of them turns on the
relation between an attribution and its grounds, so none is answerable from a table of
scores. Two are worth reproducing here.

The first asks which artifacts are attributed on detector evidence alone, in a stratum
whose measured false positive rate exceeds five per cent, with no declared evidence to
corroborate them. These are the attributions a prevalence estimate rests on, and the
query returns them ranked by genre and detector. The second asks for claims whose
supporting evidence comes from a stratum where $J$ is at or below zero. It returns the
claims that carry no information whatever their score, which is a statement no score
threshold can express and which the model makes queryable.

A third question runs in the opposite direction: artifacts an agent is declared to have
written that no detector flagged. That is the recall gap of the detector stack measured
against ground truth, and in the sample graph it is largest on diffs and smallest on
pull request bodies, which matches Table~\ref{tab:rq1}.

The vocabulary, the graph and the queries are in the replication package. The
contribution is the design and its demonstration rather than the graph itself, which is
built over a sample and serves no purpose as a dataset. We do not claim adoption: the
model is worth its length only if it makes the error rates travel with the claims, and
whether that happens depends on the platforms that issue them.

\section{Discussion}
\label{sec:discussion}

\subsection{Implications for the prevalence literature}
\label{sec:discussion:prevalence}

The measured errors do not resolve into a single multiplier that could be applied to
published figures, and the reason is itself the finding: the error depends on the
stratum. A study that scores long pull request bodies is in a different position from
one that scores commit messages, and the same detector moves from below chance on one
genre to far below chance on another.

The direction of the effect differs by channel. \citet{ji2026exploratory} apply the
zero-shot stack at thresholds imported from other domains. In our corpus those
instruments are below chance in every genre, so their flagged set is not a noisy version
of the AI-written set; it is close to unrelated to it, and the share of artifacts they
flag is closer to a property of the threshold than of the population. The keyword scans
of \citet{mao2026measurement} and \citet{almujahid2026aiusage} fail in the opposite
direction. Their false positive rate in our negative set is zero, so what they count is
real, and their recall on agent-authored text is under four per cent, so what they count
is a small and unrepresentative slice. Their numbers are lower bounds whose distance
from the truth is unknown, which is how they should be read.

Our result and that of \citet{khosravani2026census} bound the problem from opposite
sides. They show that declaration channels undercount, with bot-account lookup
recovering a small fraction of the commits their multi-method approach finds. We show
that content-based detection, offered as the remedy for exactly that gap, does not
recover the difference in this domain. A defensible prevalence figure would therefore
need a lower bound from declaration and an upper bound from corrected detection, and
Table~\ref{tab:youden} says the upper bound is not currently available. Reporting one
number from either channel alone overstates what is known.

\subsection{Implications for platform-level attribution}
\label{sec:discussion:platform}

Declared attribution works where it exists. \AidevTrailerRate{} of commits in AIDev
carry a \texttt{Co-authored-by} trailer, and the agent-account channel is exact by
construction: no inference, no threshold, no error rate to estimate. It is also cheap,
since the platform already knows what opened the pull request.

The contrast with the detection results is the argument for treating that metadata as
infrastructure rather than as a convention. Every artifact whose authorship is declared
at the moment of creation is one that no detector has to guess at later, and the
guessing, as measured here, does not work. The recommendation is narrow and
actionable: preserve agent attribution through merges, squashes and rebases, where it is
currently lost; and prefer declaration over inference wherever both are available. This
does not solve undeclared use, and nothing in this paper suggests it would.

\subsection{Implications for research practice}
\label{sec:discussion:practice}

Table~\ref{tab:checklist} turns the results into requirements for studies that estimate
how much of a corpus is AI-written. Each row names the result behind it, so none of them
rests on our judgement alone.

\begin{table}[t]
\caption{What each result asks of a prevalence study.}
\label{tab:checklist}
\small
\begin{tabular}{@{}p{0.44\textwidth}p{0.46\textwidth}@{}}
\toprule
Requirement & Because \\
\midrule
Validate on the genre being counted before estimating &
Thresholds do not transfer across domains, and every zero-shot cell in
Table~\ref{tab:zeroshot} is below chance \\
\addlinespace
Report per genre, never pooled &
Error rates differ by genre for every instrument (Tables~\ref{tab:zeroshot}
and~\ref{tab:recall}) \\
\addlinespace
Separate pre-language-model automation and report its rate &
An instructed model flags it \JudgeBotRatioMin{}--\JudgeBotRatioMax{} times more often
than human text of the same period (Table~\ref{tab:automation}) \\
\addlinespace
State the operating point with the sensitivity and specificity measured at it &
A share of artifacts flagged is uninterpretable without both \\
\addlinespace
Correct the estimate, and report when the correction is undefined &
Most instrument and genre pairs admit no correction at all
(Table~\ref{tab:estimable}) \\
\bottomrule
\end{tabular}
\end{table}

The last is the one most likely to be resisted, because it converts a headline number
into a statement that no number is available. That is nonetheless what the measurement
supports for most of the instruments in current use.

\subsection{Ethical considerations}
\label{sec:discussion:ethics}

Detector output is already used in governance: contribution policies, review triage and
moderation decisions. The error rates measured here are not evenly distributed. Rates
are highest on short artifacts and on text produced by automation, and
\citet{liang2023biased} report elevated false positives for writing by non-native
English speakers, a population heavily represented among open-source contributors. We
did not measure the linguistic dimension directly, and say so in
Section~\ref{sec:threats}; the design would require repository-level language proxies
and the claim would still be about repositories rather than people.

What the measurement does support is narrower and sufficient. An instrument scoring
below chance in every genre, used as a gate, rejects contributions at a rate unrelated
to whether they were AI-written, and the cost falls on whoever writes text the detector
finds unusual. Deploying these tools for gatekeeping is not supported by their measured
performance in this domain.

\section{Threats to validity}
\label{sec:threats}

\subsection{Construct validity}
\label{sec:threats:construct}

\paragraph{Authorship is not binary.} The strongest threat to this study is that its
outcome variable does not exist in the form the measurement assumes. Real contributions
are frequently neither purely human nor purely machine: agent output is edited by
humans before merge, and human code is refactored by agents. Treating authorship as a
binary label imposes a dichotomy on a continuum. We mitigate this in three ways. Where
possible we use the evidence-typed representation of Section~\ref{sec:provenance}
instead of a binary label; we report a sensitivity analysis excluding heavily edited
pull requests, using AIDev's review metadata; and we frame all claims as concerning
\emph{declared} agent authorship rather than AI involvement in general.

\paragraph{The unobservable stratum.} Our positive set captures agentic contributions
that platforms declare. It cannot capture a developer who pastes generated code from a
chat interface into their own editor, which is plausibly the most common form of AI
assistance in practice. This study measures detector validity on the observable
stratum, and is explicit that the unobservable one is out of scope. The N1--N2
comparison in Section~\ref{sec:groundtruth:n2} offers indirect evidence about its
magnitude under stated assumptions, but this is a bound, not a measurement.

\subsection{Internal validity}
\label{sec:threats:internal}

\paragraph{Negative-set purity.} N1's purity rests on chronology plus the requirement
that artifacts were not modified after the cutoff. Rewritten repository history could in
principle carry post-cutoff text under pre-cutoff timestamps; we report the prevalence
of detectable history rewriting in the sample. The cross-check against archived event
streams described in Section~\ref{sec:groundtruth:n1} tests the timestamp rule directly
rather than assuming it.

\paragraph{Quiescence bias.} Requiring that artifacts were never modified after the
cutoff selects for those that attracted no later activity. Such artifacts may be
systematically shorter, less discussed or from less active projects, which could bias
error-rate estimates in either direction. \RESULT{Report the observable differences
between eligible and ineligible pre-cutoff artifacts on length, project activity and
genre, so readers can judge the direction of the bias.}

\paragraph{Repository eligibility.} Only repositories predating the cutoff support
matched negatives, excluding a substantial minority of the frame
(Section~\ref{sec:groundtruth:n1}). Excluded repositories are newer and may differ in
agent adoption; the generalisation check on positives from excluded repositories tests
whether this matters.

\paragraph{Detector version drift.} Detector implementations and their underlying
scoring models change. All versions are pinned and reported, and the replication package
fixes them. This is also why no commercial detector is included: its behaviour cannot be
pinned, so a result obtained from it cannot be reproduced.

\subsection{External validity}
\label{sec:threats:external}

\paragraph{The seven-month window.} All positives fall between \AidevWindowStart{} and
\AidevWindowEnd{}. Agent platforms and their underlying models change rapidly, and
detector performance against a 2025 model generation may not transfer to later ones.
This bounds the shelf life of the specific error rates while leaving the methodological
contribution --- the counterfactual design and the correction procedure --- applicable
to whatever generation follows.

\paragraph{Agent coverage.} AIDev covers agent platforms with public GitHub footprints.
Agents operating under human accounts, and self-hosted or enterprise deployments, are
absent. Combined with the \AidevTopAgentShare{} share of the largest agent, this makes
per-agent reporting essential rather than optional.

\paragraph{Repository population.} The curated frame consists of comparatively popular
repositories, with a median of \RepoStarsMedian{} stars. Detector behaviour on small,
inactive or personal projects is not established by this study.

\subsection{Conclusion validity}
\label{sec:threats:conclusion}

The stratified analysis involves many comparisons across genres, languages, agents and
length bins. We report effect sizes with confidence intervals throughout and avoid
null-hypothesis significance language, in keeping with the reporting conventions of this
journal. Where a stratum is too thin to support a precise estimate, we report the
interval rather than the point and say so explicitly.

\section{Conclusion}
\label{sec:conclusion}

\RESULT{Restate the inversion and what measuring it produced: the measured error
structure, the corrected prevalence intervals, and the provenance model. Keep to the
evidence --- the paper's authority comes from having measured rather than argued.}

\RESULT{Close on the constructive point. The instruments are correctable and the
corrections are cheap: validate in domain, report per stratum, exclude pre-LLM
automation, and prefer declared attribution to inferred attribution wherever the choice
exists. Nothing in the finding argues that AI-generated content in open source is
unimportant --- only that the field should be able to state how well it can measure it.}


% Acknowledgements deferred to the final version.

\section*{Conflict of interest}
The author declares that he has no conflict of interest.

\section*{Data availability}
All code, the frozen corpus specification, detector configurations with pinned model
revisions, the provenance vocabulary and the SPARQL competency queries are at
\url{https://github.com/sebneu/ai-authorship-validity}, with an archived release at
\RESULT{Zenodo DOI}. Repository text is not redistributed, and neither is the AIDev
dataset \citep{li2025aidev}; the corpus is released as the specification and code needed
to rebuild it, with the exact AIDev revision pinned so that every count reported here
can be reproduced. Section~\ref{sec:groundtruth:release} describes what the release
contains.

\bibliographystyle{spbasic}
\bibliography{refs}

\end{document}
